\documentclass[11pt,reqno]{amsart}

\usepackage[T1]{fontenc}
\usepackage[utf8]{inputenc}
\usepackage{amsmath,amssymb,amsthm,mathtools}
\usepackage{booktabs}
\usepackage{float}
\usepackage{geometry}
\usepackage{microtype}
\usepackage[colorlinks=true,allcolors=blue]{hyperref}
\usepackage[nameinlink,noabbrev]{cleveref}

\geometry{margin=1in}
\numberwithin{equation}{section}

\newtheorem{theorem}{Theorem}[section]
\newtheorem{lemma}[theorem]{Lemma}
\newtheorem{proposition}[theorem]{Proposition}
\newtheorem{corollary}[theorem]{Corollary}
\theoremstyle{definition}
\newtheorem{definition}[theorem]{Definition}
\newtheorem{example}[theorem]{Example}
\theoremstyle{remark}
\newtheorem{remark}[theorem]{Remark}

\newcommand{\Rk}{\mathcal{R}_k}
\newcommand{\Sk}{\mathcal{S}_k}
\newcommand{\bk}{b_k}
\newcommand{\lk}{\ell_k}

% Edit author-facing metadata here. The manuscript body remains in main.tex
% and the section files.
\newcommand{\papertitle}{An Explicit Bijection for the \texorpdfstring{$k$}{k}-Block Index}
\newcommand{\paperauthor}{Anonymous Contributor}
\newcommand{\paperdisclosure}{This manuscript was generated by GPT and released by an anonymous repository maintainer for open technical review. It has not received independent subject-matter review.}

\title[The \texorpdfstring{$k$}{k}-block index]{\papertitle}
\author{\paperauthor}
\thanks{\paperdisclosure}
\date{July 22, 2026}
\subjclass[2020]{05A17}
\keywords{integer partitions, block index, bijection, \texorpdfstring{$k$}{k}-regular partitions}

\hypersetup{
  pdftitle={An Explicit Bijection for the k-Block Index},
  pdfauthor={Anonymous Contributor},
  pdfsubject={A constructive bijection between k-regular and k-strict partitions},
  pdfkeywords={integer partitions, block index, bijection, k-regular partitions}
}


\begin{document}

\begin{abstract}
For an integer $k\geq2$, Wang and Zhang defined a $k$-block index on
$k$-regular partitions and stated an identity equating its joint distribution
with the length to the joint distribution of the $k$-alternating length and
the length on $k$-strict partitions. A direct bijection is given. The map
iterates a simultaneous transfer which lowers the largest part of every
nonterminal block by $k$ and deletes one part from the terminal block. The
successive block indices form a partition whose adjacent differences are less
than $k$. Its conjugate is $k$-strict. An explicit reverse block transfer
establishes bijectivity and gives the required preservation of size and both
statistics.
\end{abstract}

\maketitle

\section{Introduction}

For a partition $\lambda$, write $|\lambda|$ for its size and
$\ell(\lambda)$ for its length. Wang and Zhang introduced the $k$-block index
for $k$-regular partitions in their study of refinements of Euler's partition
theorem \cite{WangZhang2026}. Their definition extends the block-index setting
considered in \cite{LiWangXu2026}. In Section~7.2 of \cite{WangZhang2026},
they stated the following identity and left its proof to the reader:
\begin{equation}\label{eq:announced-identity}
 \sum_{\lambda\in\Rk}x^{\ell(\lambda)}y^{\bk(\lambda)}q^{|\lambda|}
 =
 \sum_{\pi\in\Sk}x^{\lk(\pi)}y^{\ell(\pi)}q^{|\pi|}.
\end{equation}
Here $\Rk$, $\Sk$, $\bk$, and $\lk$ are defined in
\Cref{sec:definitions}. The purpose of this paper is to give a direct
bijective proof of \eqref{eq:announced-identity}.

The construction is local. A simultaneous transfer acts on every block of a
$k$-regular partition. Repeating the transfer records a decreasing sequence
of block indices. The central point is that consecutive entries differ by at
most $k-1$. Consequently, conjugation of the sequence produces a $k$-strict
partition. The inverse is obtained by a reverse transfer that groups parts
from the bottom upward. The proof also identifies the two statistics without
recourse to a product expansion.

\section{Definitions and statement of the result}\label{sec:definitions}

A partition is a finite weakly decreasing sequence of positive integers. Its
parts are understood with multiplicity. Fix an integer $k\geq2$. A partition
is \emph{$k$-regular} if none of its parts is divisible by $k$, and it is
\emph{$k$-strict} if every part occurs fewer than $k$ times. Let $\Rk$ and
$\Sk$ denote the corresponding sets, including the empty partition.

Let $\lambda\in\Rk$ be nonempty. Its \emph{$k$-blocks} are obtained by
working from the largest unused part downward. If $a$ is the largest part not
yet assigned, the next block contains all remaining parts with values in the
interval $[a-k,a]$. The block has weight $k$ when $a>k$, and weight $a$ when
$a\leq k$. The \emph{$k$-block index} is the sum of the block weights and is
denoted by $\bk(\lambda)$; set $\bk(\varnothing)=0$. Since the only possible
block with largest part at most $k$ is the final block, there is at most one
block whose weight is less than $k$.

For $\pi=(\pi_1,\pi_2,\ldots)\in\Sk$, extend the sequence by
$\pi_j=0$ past its length and define
\begin{equation}\label{eq:k-alternating-length}
 \lk(\pi)=\sum_{j\geq0}\bigl(\pi_{jk+1}-\pi_{(j+1)k}\bigr).
\end{equation}
The sum is finite. This is the statistic used in \cite{WangZhang2026}.

\begin{theorem}\label{thm:main}
For every integer $k\geq2$, there is a bijection $\Phi_k:\Rk\to\Sk$ such
that, for every $\lambda\in\Rk$,
\begin{equation}\label{eq:statistic-preservation}
 |\Phi_k(\lambda)|=|\lambda|,\qquad
 \ell\bigl(\Phi_k(\lambda)\bigr)=\bk(\lambda),\qquad
 \lk\bigl(\Phi_k(\lambda)\bigr)=\ell(\lambda).
\end{equation}
Consequently, \eqref{eq:announced-identity} holds.
\end{theorem}

\section{The block transfer}\label{sec:transfer}

For a partition $\nu$, let
\[
 \operatorname{supp}(\nu)=\{a\geq1:m_a(\nu)>0\},
\]
where $m_a(\nu)$ is the multiplicity of the part $a$. The transfer below
changes multiplicities, while the block decompositions are determined by this
finite support set.

\begin{lemma}[Greedy interval lemma]\label{lem:greedy-interval}
Let $D$ be a finite set of integers and let $k\geq1$. Partition $D$ by the
following two greedy procedures:
\begin{enumerate}
\item start with the least unused element $b$, take all unused elements in
$[b,b+k]$, and continue upward;
\item start with the greatest unused element $a$, take all unused elements in
$[a-k,a]$, and continue downward.
\end{enumerate}
The two procedures produce the same number of blocks.
\end{lemma}

\begin{proof}
In the first procedure, record the least element selected at each step. These
recorded elements differ pairwise by more than $k$, and every subset of $D$
whose distinct elements differ pairwise by more than $k$ contains at most one
element in each greedy block. Thus the number of blocks in the first
procedure is the maximum cardinality of a $k$-separated subset of $D$.
The greatest-first procedure has the same property, using its greatest
elements as the recorded separated subset. Hence the two block counts agree.
\end{proof}

Let $\lambda\in\Rk$ be nonempty. Define $\tau(\lambda)$ by acting
simultaneously on its $k$-blocks. From every block with largest part $a>k$,
replace one occurrence of $a$ by $a-k$. From the final block with largest
part $a\leq k$, if it exists, delete one occurrence of $a$. The resulting
partition is again $k$-regular. The decrease in size is the sum of the block
weights:
\begin{equation}\label{eq:size-transfer}
 |\tau(\lambda)|=|\lambda|-\bk(\lambda).
\end{equation}

The reverse operation is explicit. For an integer $w>0$, let $r$ be its least
nonnegative residue modulo $k$.

\begin{lemma}[Explicit reverse transfer]\label{lem:block-transfer}
Let $\mu\in\Rk$, put $h=\bk(\mu)$, and suppose that
\begin{equation}\label{eq:admissible-weight}
 h\leq w\leq h+k-1.
\end{equation}
Construct $\sigma_w(\mu)$ by the following finite algorithm.
\begin{enumerate}
\item If $r>0$, adjoin one part $r$ to $\mu$.
\item Leave all parts of size at most $r$ fixed. On the distinct values in
$\operatorname{supp}(\mu)$ larger than $r$, work from the smallest upward. If
$b$ is the smallest unprocessed value, make a reverse block from all
unprocessed values in $[b,b+k]$, and continue with values larger than $b+k$.
\item In every reverse block, replace one occurrence of its smallest part $b$
by $b+k$.
\end{enumerate}
Then $\sigma_w(\mu)\in\Rk$ and
\begin{equation}\label{eq:transfer-inverse}
 \tau\bigl(\sigma_w(\mu)\bigr)=\mu,
 \qquad
 \bk\bigl(\sigma_w(\mu)\bigr)=w.
\end{equation}
\end{lemma}

\begin{proof}
The algorithm is deterministic and terminates because the support is finite.
Every new value has the same nonzero residue modulo $k$ as the value it
replaces, so $\sigma_w(\mu)$ is $k$-regular.

Write
\begin{equation}\label{eq:decompose-h}
 h=ak+t,\qquad 0\leq t<k.
\end{equation}
There are $a$ blocks of $\mu$ whose largest parts exceed $k$; if $t>0$, the
remaining block has largest part $t$. Let
\[
 D_r=\{d\in\operatorname{supp}(\mu):d>r\}.
\]
Removing the values at most $r$ removes the terminal block when $r\geq t$,
and leaves a nonempty part of it when $r<t$. The blocks with largest parts
greater than $k$ retain their largest parts, so the greatest-first
decomposition of $D_r$ has $a$ blocks when $r\geq t$, and $a+1$ blocks when
$r<t$. By \Cref{lem:greedy-interval}, the reverse-block algorithm in step 2
has the same number, say $q$, of blocks. Consequently $q=a$ in the first
case and $q=a+1$ in the second. The interval condition
\eqref{eq:admissible-weight} gives $w=ak+r$ in the first case and
$w=(a+1)k+r$ in the second. Therefore $qk+r=w$.

Let the reverse blocks have smallest values $b_1<\cdots<b_q$. Consecutive
smallest values satisfy $b_{j+1}>b_j+k$. After step 3, the values originating
in the $j$th reverse block lie in $[b_j,b_j+k]$ and include $b_j+k$.
Moreover, $b_j+k<b_{j+1}$, so the top-down decomposition cannot merge two
adjacent reverse blocks. It therefore recovers these intervals, in reverse
order, as blocks with largest values $b_q+k,\ldots,b_1+k$. If $r>0$, all
fixed values are at most $r$ and the adjoined value $r$ is present. When
$q>0$, they lie below $b_1$ and form the final block; when $q=0$, they form
the only block. In both cases its largest value is $r$.

Thus the block weights of $\sigma_w(\mu)$ are $q$ copies of $k$ and, when
$r>0$, one weight $r$; their sum is $qk+r=w$. Applying $\tau$ lowers one
$b_j+k$ to $b_j$ in every nonterminal block and removes the adjoined part
$r$ from the terminal block. All multiplicities are thereby restored, so
$\tau(\sigma_w(\mu))=\mu$.
\end{proof}

\begin{lemma}[Forward transfer is inverted by the algorithm]
\label{lem:converse-transfer}
Let $\lambda\in\Rk$ be nonempty, put $w=\bk(\lambda)$, and set
$\mu=\tau(\lambda)$. Then
\begin{equation}\label{eq:block-drop}
 \bk(\mu)\leq w\leq\bk(\mu)+k-1,
\end{equation}
and $\lambda=\sigma_w(\mu)$.
\end{lemma}

\begin{proof}
Write $w=ak+r$ with $0\leq r<k$. The $a$ nonterminal blocks of $\lambda$ have
largest parts $A_1>\cdots>A_a>k$. If $r>0$, the final block has largest
part $r$; if $r=0$, there is no final block of weight less than $k$.

For $1\leq j\leq a$, put $c_j=A_j-k$. The image under $\tau$ of the $j$th
nonterminal block has support contained in $[c_j,c_j+k]$ and contains
$c_j$, because one occurrence of $A_j$ was moved to $A_j-k$. For
$1\leq j<a$, the next lower nonterminal block satisfies
\[
 c_{j+1}+k=A_{j+1}<A_j-k=c_j,
\]
by the defining greedy separation of the original blocks. Thus the support
intervals $[c_j,c_j+k]$ are pairwise separated. If a terminal block exists,
its largest part is $r<c_a$; after the deletion, every surviving terminal
value is at most $r$.

Consequently, the distinct values of $\mu$ larger than $r$ are partitioned by
the reverse greedy procedure into exactly the $a$ transformed nonterminal
blocks. It starts at $c_a$, recovers the support in $[c_a,c_a+k]$, and then
proceeds successively through the separated intervals up to $[c_1,c_1+k]$.
The values at most $r$ are precisely the unchanged terminal values. Step 3 of
the construction of $\sigma_w(\mu)$ therefore reverses every move made by
$\tau$, including the possible terminal deletion. Hence $\lambda=\sigma_w(\mu)$.

It remains to compare block indices. Set
\[
 D_r=\{d\in\operatorname{supp}(\mu):d>r\}.
\]
If $D_r$ is empty, then $a=0$ and all parts of $\mu$ are at most $r$.
Hence $\bk(\mu)=t$ for some $0\leq t\leq r$, where $t=0$ means that $\mu$
is empty.

Suppose that $D_r$ is nonempty. Its bottom-up decomposition has the $a$
transformed nonterminal blocks identified above, so \Cref{lem:greedy-interval}
shows that its top-down decomposition also has exactly $a$ blocks. Since $\mu$
is $k$-regular, the top of its lowest block cannot equal $k$. If that top is
greater than $k$, then all $a$ blocks have weight $k$. Values at most $r$ either
join the lowest block or form a terminal block of top $t\leq r$; in either event
\[
 \bk(\mu)=ak+t\qquad(0\leq t\leq r),
\]
with $t=0$ when there is no terminal block. If the lowest block has top less
than $k$, it is itself the terminal block and contains every positive value
of $\mu$ at most $r$. Its top $t$ satisfies $r<t<k$, while the other $a-1$
blocks have weight $k$. Thus
\[
 \bk(\mu)=(a-1)k+t\qquad(r<t<k).
\]
The first formula also covers $D_r=\varnothing$. In the two cases,
respectively,
\[
 w-\bk(\mu)=r-t,
 \qquad
 w-\bk(\mu)=k+r-t.
\]
Both differences lie in $\{0,1,\ldots,k-1\}$, which proves
\eqref{eq:block-drop}.
\end{proof}

\section{The bijection and its statistics}\label{sec:bijection}

Start with $\lambda^{(0)}=\lambda\in\Rk$ and define
\begin{equation}\label{eq:iterated-transfer}
 \lambda^{(i+1)}=\tau\bigl(\lambda^{(i)}\bigr)
 \quad\text{while }\lambda^{(i)}\ne\varnothing.
\end{equation}
The process terminates by \eqref{eq:size-transfer}. Let
\begin{equation}\label{eq:transfer-sequence}
 \mathbf{s}(\lambda)=(s_1,s_2,\ldots,s_m),
 \qquad s_i=\bk\bigl(\lambda^{(i-1)}\bigr),
\end{equation}
where $\lambda^{(m)}=\varnothing$. Define
\begin{equation}\label{eq:definition-phi}
 \Phi_k(\lambda)=\mathbf{s}(\lambda)'.
\end{equation}

\begin{proposition}\label{prop:sequence-properties}
Let $\lambda\in\Rk$ be nonempty. The sequence $\mathbf{s}(\lambda)$ is a
partition and satisfies
\begin{equation}\label{eq:gap-bound}
 0\leq s_i-s_{i+1}<k\quad(1\leq i<m),
 \qquad 0<s_m<k.
\end{equation}
Moreover, $\Phi_k(\lambda)$ is $k$-strict and satisfies the three identities
in \eqref{eq:statistic-preservation}.
\end{proposition}

\begin{proof}
Apply \Cref{lem:converse-transfer} at every step of
\eqref{eq:iterated-transfer}. It gives the first inequality in
\eqref{eq:gap-bound}. In the final nonempty partition all remaining parts lie
in a single block with top less than $k$, so $0<s_m<k$.

For a partition $\mathbf{s}$, the multiplicity of the part $i$ in its
conjugate is $s_i-s_{i+1}$, where $s_{m+1}=0$. Thus
\eqref{eq:gap-bound} implies that $\mathbf{s}'$ is $k$-strict. Telescoping
\eqref{eq:size-transfer} over the transfer sequence gives
\begin{equation}\label{eq:size-telescope}
 |\lambda|=s_1+s_2+\cdots+s_m=|\mathbf{s}'|.
\end{equation}
The length of $\mathbf{s}'$ is its largest part $s_1=\bk(\lambda)$.

It remains to identify the length statistic. A transfer changes no number of
parts except when a terminal block is present, in which case it deletes one
part. Such a block is present exactly when the corresponding $s_i$ is not a
multiple of $k$. Hence
\begin{equation}\label{eq:length-nonmultiples}
 \ell(\lambda)=\#\{i: k\nmid s_i\}.
\end{equation}
Let $\pi=\mathbf{s}'$. For every $j\geq0$,
\begin{equation}\label{eq:strip-count}
 \pi_{jk+1}-\pi_{(j+1)k}
 =\#\{i: jk<s_i<(j+1)k\}.
\end{equation}
Summing \eqref{eq:strip-count} and using \eqref{eq:k-alternating-length}
shows that its left-hand side is $\lk(\pi)$ and its right-hand side is the
quantity in \eqref{eq:length-nonmultiples}. This proves all three identities.
\end{proof}

\begin{proof}[Proof of \Cref{thm:main}]
The empty partition maps to itself. For a nonempty $\lambda$, the preceding
proposition shows that $\Phi_k$ has the required statistics. It remains to
prove that the map is a bijection on nonempty partitions. Let $\pi\in\Sk$ be
nonempty and put
$\mathbf{s}=\pi'=(s_1,\ldots,s_m)$. Since $\pi$ is $k$-strict,
\begin{equation}\label{eq:inverse-gap-bound}
 0\leq s_i-s_{i+1}<k\quad(1\leq i<m),
 \qquad 0<s_m<k.
\end{equation}
Set $\mu^{(m)}=\varnothing$. For $i=m,m-1,\ldots,1$, define recursively
\begin{equation}\label{eq:inverse-recursion}
 \mu^{(i-1)}=\sigma_{s_i}\bigl(\mu^{(i)}\bigr).
\end{equation}
Inductively, $\bk(\mu^{(i)})=s_{i+1}$, with $s_{m+1}=0$. Thus
\eqref{eq:inverse-gap-bound} is exactly the admissibility condition in
\Cref{lem:block-transfer}. That lemma gives
$\tau(\mu^{(i-1)})=\mu^{(i)}$ and
$\bk(\mu^{(i-1)})=s_i$. Therefore
$\mathbf{s}(\mu^{(0)})=\mathbf{s}$, so $\Phi_k(\mu^{(0)})=\pi$. The same
reverse recursion is forced by \Cref{lem:converse-transfer}; hence the
preimage is unique. This proves bijectivity and the theorem.
\end{proof}

\section{An example}\label{sec:example}

Take $k=4$ and $\lambda=(10,7,3,1)$. Its first block is $(10,7)$ and its
terminal block is $(3,1)$, so $\bk(\lambda)=7$. Iterating the transfer gives
the following table.

\begin{table}[H]
\centering
\caption{The transfer sequence for $k=4$ and $\lambda=(10,7,3,1)$.}
\label{tab:transfer-example}
\begin{tabular}{ccl}
\toprule
$i$ & $\lambda^{(i)}$ & $\bk(\lambda^{(i)})$ \\
\midrule
0 & $(10,7,3,1)$ & $7$ \\
1 & $(7,6,1)$ & $5$ \\
2 & $(6,3)$ & $4$ \\
3 & $(3,2)$ & $3$ \\
4 & $(2)$ & $2$ \\
\bottomrule
\end{tabular}
\end{table}

Thus $\mathbf{s}(\lambda)=(7,5,4,3,2)$ and
\begin{equation*}
 \Phi_4(\lambda)=(5,5,4,3,2,1,1).
\end{equation*}
This partition is $4$-strict. Its size is $21$, its length is $7$, and
\begin{equation*}
 \ell_4\bigl(\Phi_4(\lambda)\bigr)
 =(5-3)+(2-0)=4=\ell(\lambda),
\end{equation*}
which displays the three statistic equalities in \Cref{thm:main}.


\bibliographystyle{amsplain}
\bibliography{references}

\end{document}
